% 09_clock_power.tex -- PLL, clock tree + gating, DVFS, power domains, reset synchronizer, CDC.
% Build: tectonic 09_clock_power.tex
\documentclass[tikz,border=10pt,12pt]{standalone}
% soc_style.tex -- shared look for the SoC figures (same palette as the slide decks and the
% core register placement figure). Every figure does % soc_style.tex -- shared look for the SoC figures (same palette as the slide decks and the
% core register placement figure). Every figure does % soc_style.tex -- shared look for the SoC figures (same palette as the slide decks and the
% core register placement figure). Every figure does \input{soc_style}.
\usepackage[T1]{fontenc}
\usepackage[scaled=0.95]{helvet}
\renewcommand{\familydefault}{\sfdefault}
\usepackage{booktabs}
\usepackage{amsmath}
\usepackage{amssymb}
\usetikzlibrary{arrows.meta,calc,positioning,fit,shapes.geometric,shapes.misc}

% ---------------- palette ----------------
\definecolor{pcb}{HTML}{0B2E33}
\definecolor{pcbmid}{HTML}{1F5C61}
\definecolor{tint}{HTML}{E6F0EF}
\definecolor{copper}{HTML}{B5651D}
\definecolor{coppert}{HTML}{F6E8DC}
\definecolor{ink}{HTML}{1F2A2E}
\definecolor{muted}{HTML}{5B6B70}
\definecolor{hardf}{HTML}{DCE6F2}
\definecolor{harde}{HTML}{2F4A7A}
\definecolor{good}{HTML}{2E7D32}
\definecolor{bad}{HTML}{B3261E}
% one colour per kind of block: border (d) / fill (f)
\definecolor{cpud}{HTML}{2B5C9E} \definecolor{cpuf}{HTML}{D6E4F5}
\definecolor{gpud}{HTML}{2E7D32} \definecolor{gpuf}{HTML}{D5EDD0}
\definecolor{npud}{HTML}{0B2E33} \definecolor{npuf}{HTML}{BFE0DB}
\definecolor{memd}{HTML}{5E3A8C} \definecolor{memf}{HTML}{E4DAF3}
\definecolor{busd}{HTML}{B5651D} \definecolor{busf}{HTML}{F6E8DC}
\definecolor{prd}{HTML}{546E7A}  \definecolor{prf}{HTML}{E3EAED}
\definecolor{secd}{HTML}{B3261E} \definecolor{secf}{HTML}{F7D9D6}
\definecolor{pwrd}{HTML}{8A6800} \definecolor{pwrf}{HTML}{FFF0B8}
\definecolor{medd}{HTML}{8E2463} \definecolor{medf}{HTML}{F5D6E6}
\definecolor{skd}{HTML}{2E7D32}  \definecolor{skf}{HTML}{BFE3B4}   % skew buffers, as in the register figure

\tikzset{
  blk/.style 2 args={draw=#1, fill=#2, rounded corners=4pt, thick, align=center,
                     font=\scriptsize, inner sep=3pt},
  cpu/.style={blk={cpud}{cpuf}},  gpu/.style={blk={gpud}{gpuf}},
  npu/.style={blk={npud}{npuf}},  mem/.style={blk={memd}{memf}},
  bus/.style={blk={busd}{busf}},  per/.style={blk={prd}{prf}},
  sec/.style={blk={secd}{secf}},  pwr/.style={blk={pwrd}{pwrf}},
  med/.style={blk={medd}{medf}},
  comb/.style ={draw=copper, fill=coppert, rounded corners=2pt, thick, align=center,
                font=\scriptsize, inner sep=2pt},                       % combinational logic
  hard/.style ={draw=harde, fill=hardf, thick, dashed, align=center,
                font=\scriptsize, inner sep=2pt},                       % hard macro / off-chip
  sram/.style ={draw=memd, fill=memf, thick, align=center, font=\scriptsize,
                inner sep=2pt},                                         % SRAM array, not flip-flops
  ff/.style   ={draw=#1, fill=#1!30, line width=0.35pt},                % one flip-flop square
  ffgone/.style={draw=muted, fill=white, densely dashed, line width=0.35pt},
  panel/.style={draw=muted, rounded corners=6pt, line width=0.9pt},
  ptitle/.style={anchor=north west, font=\normalsize\bfseries, text=ink, inner sep=0pt},
  note/.style ={font=\scriptsize, text=ink, align=left, inner sep=0pt},
  lab/.style  ={font=\tiny, text=muted, inner sep=1pt},
  axi/.style  ={{Stealth[length=2.6mm]}-{Stealth[length=2.6mm]}, line width=2.2pt, copper},
  axi1/.style ={-{Stealth[length=2.6mm]}, line width=2.2pt, copper},
  apb/.style  ={{Stealth[length=1.8mm]}-{Stealth[length=1.8mm]}, line width=1pt, copper},
  apb1/.style ={-{Stealth[length=1.8mm]}, line width=1pt, copper},
  dat/.style  ={-{Stealth[length=1.8mm]}, thick, pcbmid},
  ctl/.style  ={-{Stealth[length=1.8mm]}, densely dashed, thick, muted},
  irq/.style  ={-{Stealth[length=2mm]}, densely dotted, line width=1.3pt, bad},
  num/.style  ={circle, fill=copper, text=white, font=\scriptsize\bfseries, inner sep=1pt,
                minimum size=0.46cm},
  ours/.style ={draw=bad, line width=1.2pt, densely dashed, rounded corners=4pt},
}

% ---------------- helpers ----------------
\newcommand{\rn}[1]{\texttt{#1}}
% \ffrow{x}{y}{n}{cell}{style}: n flip-flops in a row, bottom-left corner at (x,y)
\newcommand{\ffrow}[5]{\foreach \q in {1,...,#3}{\draw[#5] ({#1+(\q-1)*#4},{#2}) rectangle ++(#4,#4);}}
% \ffcol{x}{y}{n}{cell}{style}: n flip-flops in a column, top-left corner at (x,y)
\newcommand{\ffcol}[5]{\foreach \q in {1,...,#3}{\draw[#5] ({#1},{#2-\q*#4}) rectangle ++(#4,#4);}}
% \ffgrid{x}{ytop}{n}{per row}{cell}{style}: n flip-flops, filled row by row from the top-left
\newcommand{\ffgrid}[6]{\pgfmathtruncatemacro{\gl}{#3-1}%
  \foreach \b in {0,...,\gl}{\pgfmathtruncatemacro{\gc}{mod(\b,#4)}\pgfmathtruncatemacro{\gr}{div(\b,#4)}%
    \draw[#6] ({#1+\gc*#5},{#2-(\gr+1)*#5}) rectangle ++(#5,#5);}}

% every figure: same page size, title, footer
\newcommand{\canvas}{\useasboundingbox (-0.3,-0.2) rectangle (38.6,30.4);}
\newcommand{\figtitle}[3]{%
  \node[num, minimum size=0.9cm, font=\Large\bfseries] at (0.45,29.55) {#1};
  \node[anchor=west, font=\LARGE\bfseries, text=pcb] at (1.15,29.55) {#2};
  \node[anchor=west, font=\small, text=ink] at (0,28.7) {#3};}
\newcommand{\figfoot}[1]{\node[anchor=west, font=\scriptsize, text=muted] at (0,0.25) {#1};}
% \panel{x0}{y0}{x1}{y1}{title}
\newcommand{\panel}[5]{\draw[panel] (#1,#2) rectangle (#3,#4);
  \node[ptitle] at ({#1+0.3},{#4-0.25}) {#5};}

% ---------------- timing diagrams: 1 cycle = \cw cm, low at y, high at y+0.45 ----------------
\newcommand{\cw}{1.0}
\newcommand{\wclk}[3]{% x0 y ncycles
  \foreach \k in {1,...,#3}{\draw[pcb, thick] ({#1+(\k-1)*\cw},#2) -- ++(0,0.45) -- ++(0.5*\cw,0)
     -- ++(0,-0.45) -- ++(0.5*\cw,0);}}
\newcommand{\wsig}[3]{% x0 y {level per cycle, e.g. 0,1,1,0}
  \gdef\wprev{x}%
  \foreach \v [count=\k from 0] in {#3}{%
    \draw[pcb, thick] ({#1+\k*\cw},{#2+0.45*\v}) -- ++(\cw,0);
    \if x\wprev\else\ifnum\wprev=\v\else\draw[pcb, thick] ({#1+\k*\cw},#2) -- ++(0,0.45);\fi\fi
    \global\let\wprev\v}}
\newcommand{\wbus}[5]{% x0 y c0 c1 text: one value held from cycle c0 to cycle c1
  \draw[pcbmid, thick, fill=white] ({#1+(#3)*\cw},{#2+0.225}) -- ({#1+(#3)*\cw+0.08},{#2+0.45})
    -- ({#1+(#4)*\cw-0.08},{#2+0.45}) -- ({#1+(#4)*\cw},{#2+0.225}) -- ({#1+(#4)*\cw-0.08},#2)
    -- ({#1+(#3)*\cw+0.08},#2) -- cycle;
  \node[font=\tiny\ttfamily] at ({#1+((#3)+(#4))*\cw/2},{#2+0.225}) {#5};}
\newcommand{\widle}[4]{% x0 y c0 c1: bus idle (no value)
  \draw[muted, thick] ({#1+(#3)*\cw},{#2+0.225}) -- ({#1+(#4)*\cw},{#2+0.225});}
\newcommand{\wname}[3]{\node[anchor=east, font=\scriptsize\ttfamily] at ({#1-0.15},{#2+0.225}) {#3};}
.
\usepackage[T1]{fontenc}
\usepackage[scaled=0.95]{helvet}
\renewcommand{\familydefault}{\sfdefault}
\usepackage{booktabs}
\usepackage{amsmath}
\usepackage{amssymb}
\usetikzlibrary{arrows.meta,calc,positioning,fit,shapes.geometric,shapes.misc}

% ---------------- palette ----------------
\definecolor{pcb}{HTML}{0B2E33}
\definecolor{pcbmid}{HTML}{1F5C61}
\definecolor{tint}{HTML}{E6F0EF}
\definecolor{copper}{HTML}{B5651D}
\definecolor{coppert}{HTML}{F6E8DC}
\definecolor{ink}{HTML}{1F2A2E}
\definecolor{muted}{HTML}{5B6B70}
\definecolor{hardf}{HTML}{DCE6F2}
\definecolor{harde}{HTML}{2F4A7A}
\definecolor{good}{HTML}{2E7D32}
\definecolor{bad}{HTML}{B3261E}
% one colour per kind of block: border (d) / fill (f)
\definecolor{cpud}{HTML}{2B5C9E} \definecolor{cpuf}{HTML}{D6E4F5}
\definecolor{gpud}{HTML}{2E7D32} \definecolor{gpuf}{HTML}{D5EDD0}
\definecolor{npud}{HTML}{0B2E33} \definecolor{npuf}{HTML}{BFE0DB}
\definecolor{memd}{HTML}{5E3A8C} \definecolor{memf}{HTML}{E4DAF3}
\definecolor{busd}{HTML}{B5651D} \definecolor{busf}{HTML}{F6E8DC}
\definecolor{prd}{HTML}{546E7A}  \definecolor{prf}{HTML}{E3EAED}
\definecolor{secd}{HTML}{B3261E} \definecolor{secf}{HTML}{F7D9D6}
\definecolor{pwrd}{HTML}{8A6800} \definecolor{pwrf}{HTML}{FFF0B8}
\definecolor{medd}{HTML}{8E2463} \definecolor{medf}{HTML}{F5D6E6}
\definecolor{skd}{HTML}{2E7D32}  \definecolor{skf}{HTML}{BFE3B4}   % skew buffers, as in the register figure

\tikzset{
  blk/.style 2 args={draw=#1, fill=#2, rounded corners=4pt, thick, align=center,
                     font=\scriptsize, inner sep=3pt},
  cpu/.style={blk={cpud}{cpuf}},  gpu/.style={blk={gpud}{gpuf}},
  npu/.style={blk={npud}{npuf}},  mem/.style={blk={memd}{memf}},
  bus/.style={blk={busd}{busf}},  per/.style={blk={prd}{prf}},
  sec/.style={blk={secd}{secf}},  pwr/.style={blk={pwrd}{pwrf}},
  med/.style={blk={medd}{medf}},
  comb/.style ={draw=copper, fill=coppert, rounded corners=2pt, thick, align=center,
                font=\scriptsize, inner sep=2pt},                       % combinational logic
  hard/.style ={draw=harde, fill=hardf, thick, dashed, align=center,
                font=\scriptsize, inner sep=2pt},                       % hard macro / off-chip
  sram/.style ={draw=memd, fill=memf, thick, align=center, font=\scriptsize,
                inner sep=2pt},                                         % SRAM array, not flip-flops
  ff/.style   ={draw=#1, fill=#1!30, line width=0.35pt},                % one flip-flop square
  ffgone/.style={draw=muted, fill=white, densely dashed, line width=0.35pt},
  panel/.style={draw=muted, rounded corners=6pt, line width=0.9pt},
  ptitle/.style={anchor=north west, font=\normalsize\bfseries, text=ink, inner sep=0pt},
  note/.style ={font=\scriptsize, text=ink, align=left, inner sep=0pt},
  lab/.style  ={font=\tiny, text=muted, inner sep=1pt},
  axi/.style  ={{Stealth[length=2.6mm]}-{Stealth[length=2.6mm]}, line width=2.2pt, copper},
  axi1/.style ={-{Stealth[length=2.6mm]}, line width=2.2pt, copper},
  apb/.style  ={{Stealth[length=1.8mm]}-{Stealth[length=1.8mm]}, line width=1pt, copper},
  apb1/.style ={-{Stealth[length=1.8mm]}, line width=1pt, copper},
  dat/.style  ={-{Stealth[length=1.8mm]}, thick, pcbmid},
  ctl/.style  ={-{Stealth[length=1.8mm]}, densely dashed, thick, muted},
  irq/.style  ={-{Stealth[length=2mm]}, densely dotted, line width=1.3pt, bad},
  num/.style  ={circle, fill=copper, text=white, font=\scriptsize\bfseries, inner sep=1pt,
                minimum size=0.46cm},
  ours/.style ={draw=bad, line width=1.2pt, densely dashed, rounded corners=4pt},
}

% ---------------- helpers ----------------
\newcommand{\rn}[1]{\texttt{#1}}
% \ffrow{x}{y}{n}{cell}{style}: n flip-flops in a row, bottom-left corner at (x,y)
\newcommand{\ffrow}[5]{\foreach \q in {1,...,#3}{\draw[#5] ({#1+(\q-1)*#4},{#2}) rectangle ++(#4,#4);}}
% \ffcol{x}{y}{n}{cell}{style}: n flip-flops in a column, top-left corner at (x,y)
\newcommand{\ffcol}[5]{\foreach \q in {1,...,#3}{\draw[#5] ({#1},{#2-\q*#4}) rectangle ++(#4,#4);}}
% \ffgrid{x}{ytop}{n}{per row}{cell}{style}: n flip-flops, filled row by row from the top-left
\newcommand{\ffgrid}[6]{\pgfmathtruncatemacro{\gl}{#3-1}%
  \foreach \b in {0,...,\gl}{\pgfmathtruncatemacro{\gc}{mod(\b,#4)}\pgfmathtruncatemacro{\gr}{div(\b,#4)}%
    \draw[#6] ({#1+\gc*#5},{#2-(\gr+1)*#5}) rectangle ++(#5,#5);}}

% every figure: same page size, title, footer
\newcommand{\canvas}{\useasboundingbox (-0.3,-0.2) rectangle (38.6,30.4);}
\newcommand{\figtitle}[3]{%
  \node[num, minimum size=0.9cm, font=\Large\bfseries] at (0.45,29.55) {#1};
  \node[anchor=west, font=\LARGE\bfseries, text=pcb] at (1.15,29.55) {#2};
  \node[anchor=west, font=\small, text=ink] at (0,28.7) {#3};}
\newcommand{\figfoot}[1]{\node[anchor=west, font=\scriptsize, text=muted] at (0,0.25) {#1};}
% \panel{x0}{y0}{x1}{y1}{title}
\newcommand{\panel}[5]{\draw[panel] (#1,#2) rectangle (#3,#4);
  \node[ptitle] at ({#1+0.3},{#4-0.25}) {#5};}

% ---------------- timing diagrams: 1 cycle = \cw cm, low at y, high at y+0.45 ----------------
\newcommand{\cw}{1.0}
\newcommand{\wclk}[3]{% x0 y ncycles
  \foreach \k in {1,...,#3}{\draw[pcb, thick] ({#1+(\k-1)*\cw},#2) -- ++(0,0.45) -- ++(0.5*\cw,0)
     -- ++(0,-0.45) -- ++(0.5*\cw,0);}}
\newcommand{\wsig}[3]{% x0 y {level per cycle, e.g. 0,1,1,0}
  \gdef\wprev{x}%
  \foreach \v [count=\k from 0] in {#3}{%
    \draw[pcb, thick] ({#1+\k*\cw},{#2+0.45*\v}) -- ++(\cw,0);
    \if x\wprev\else\ifnum\wprev=\v\else\draw[pcb, thick] ({#1+\k*\cw},#2) -- ++(0,0.45);\fi\fi
    \global\let\wprev\v}}
\newcommand{\wbus}[5]{% x0 y c0 c1 text: one value held from cycle c0 to cycle c1
  \draw[pcbmid, thick, fill=white] ({#1+(#3)*\cw},{#2+0.225}) -- ({#1+(#3)*\cw+0.08},{#2+0.45})
    -- ({#1+(#4)*\cw-0.08},{#2+0.45}) -- ({#1+(#4)*\cw},{#2+0.225}) -- ({#1+(#4)*\cw-0.08},#2)
    -- ({#1+(#3)*\cw+0.08},#2) -- cycle;
  \node[font=\tiny\ttfamily] at ({#1+((#3)+(#4))*\cw/2},{#2+0.225}) {#5};}
\newcommand{\widle}[4]{% x0 y c0 c1: bus idle (no value)
  \draw[muted, thick] ({#1+(#3)*\cw},{#2+0.225}) -- ({#1+(#4)*\cw},{#2+0.225});}
\newcommand{\wname}[3]{\node[anchor=east, font=\scriptsize\ttfamily] at ({#1-0.15},{#2+0.225}) {#3};}
.
\usepackage[T1]{fontenc}
\usepackage[scaled=0.95]{helvet}
\renewcommand{\familydefault}{\sfdefault}
\usepackage{booktabs}
\usepackage{amsmath}
\usepackage{amssymb}
\usetikzlibrary{arrows.meta,calc,positioning,fit,shapes.geometric,shapes.misc}

% ---------------- palette ----------------
\definecolor{pcb}{HTML}{0B2E33}
\definecolor{pcbmid}{HTML}{1F5C61}
\definecolor{tint}{HTML}{E6F0EF}
\definecolor{copper}{HTML}{B5651D}
\definecolor{coppert}{HTML}{F6E8DC}
\definecolor{ink}{HTML}{1F2A2E}
\definecolor{muted}{HTML}{5B6B70}
\definecolor{hardf}{HTML}{DCE6F2}
\definecolor{harde}{HTML}{2F4A7A}
\definecolor{good}{HTML}{2E7D32}
\definecolor{bad}{HTML}{B3261E}
% one colour per kind of block: border (d) / fill (f)
\definecolor{cpud}{HTML}{2B5C9E} \definecolor{cpuf}{HTML}{D6E4F5}
\definecolor{gpud}{HTML}{2E7D32} \definecolor{gpuf}{HTML}{D5EDD0}
\definecolor{npud}{HTML}{0B2E33} \definecolor{npuf}{HTML}{BFE0DB}
\definecolor{memd}{HTML}{5E3A8C} \definecolor{memf}{HTML}{E4DAF3}
\definecolor{busd}{HTML}{B5651D} \definecolor{busf}{HTML}{F6E8DC}
\definecolor{prd}{HTML}{546E7A}  \definecolor{prf}{HTML}{E3EAED}
\definecolor{secd}{HTML}{B3261E} \definecolor{secf}{HTML}{F7D9D6}
\definecolor{pwrd}{HTML}{8A6800} \definecolor{pwrf}{HTML}{FFF0B8}
\definecolor{medd}{HTML}{8E2463} \definecolor{medf}{HTML}{F5D6E6}
\definecolor{skd}{HTML}{2E7D32}  \definecolor{skf}{HTML}{BFE3B4}   % skew buffers, as in the register figure

\tikzset{
  blk/.style 2 args={draw=#1, fill=#2, rounded corners=4pt, thick, align=center,
                     font=\scriptsize, inner sep=3pt},
  cpu/.style={blk={cpud}{cpuf}},  gpu/.style={blk={gpud}{gpuf}},
  npu/.style={blk={npud}{npuf}},  mem/.style={blk={memd}{memf}},
  bus/.style={blk={busd}{busf}},  per/.style={blk={prd}{prf}},
  sec/.style={blk={secd}{secf}},  pwr/.style={blk={pwrd}{pwrf}},
  med/.style={blk={medd}{medf}},
  comb/.style ={draw=copper, fill=coppert, rounded corners=2pt, thick, align=center,
                font=\scriptsize, inner sep=2pt},                       % combinational logic
  hard/.style ={draw=harde, fill=hardf, thick, dashed, align=center,
                font=\scriptsize, inner sep=2pt},                       % hard macro / off-chip
  sram/.style ={draw=memd, fill=memf, thick, align=center, font=\scriptsize,
                inner sep=2pt},                                         % SRAM array, not flip-flops
  ff/.style   ={draw=#1, fill=#1!30, line width=0.35pt},                % one flip-flop square
  ffgone/.style={draw=muted, fill=white, densely dashed, line width=0.35pt},
  panel/.style={draw=muted, rounded corners=6pt, line width=0.9pt},
  ptitle/.style={anchor=north west, font=\normalsize\bfseries, text=ink, inner sep=0pt},
  note/.style ={font=\scriptsize, text=ink, align=left, inner sep=0pt},
  lab/.style  ={font=\tiny, text=muted, inner sep=1pt},
  axi/.style  ={{Stealth[length=2.6mm]}-{Stealth[length=2.6mm]}, line width=2.2pt, copper},
  axi1/.style ={-{Stealth[length=2.6mm]}, line width=2.2pt, copper},
  apb/.style  ={{Stealth[length=1.8mm]}-{Stealth[length=1.8mm]}, line width=1pt, copper},
  apb1/.style ={-{Stealth[length=1.8mm]}, line width=1pt, copper},
  dat/.style  ={-{Stealth[length=1.8mm]}, thick, pcbmid},
  ctl/.style  ={-{Stealth[length=1.8mm]}, densely dashed, thick, muted},
  irq/.style  ={-{Stealth[length=2mm]}, densely dotted, line width=1.3pt, bad},
  num/.style  ={circle, fill=copper, text=white, font=\scriptsize\bfseries, inner sep=1pt,
                minimum size=0.46cm},
  ours/.style ={draw=bad, line width=1.2pt, densely dashed, rounded corners=4pt},
}

% ---------------- helpers ----------------
\newcommand{\rn}[1]{\texttt{#1}}
% \ffrow{x}{y}{n}{cell}{style}: n flip-flops in a row, bottom-left corner at (x,y)
\newcommand{\ffrow}[5]{\foreach \q in {1,...,#3}{\draw[#5] ({#1+(\q-1)*#4},{#2}) rectangle ++(#4,#4);}}
% \ffcol{x}{y}{n}{cell}{style}: n flip-flops in a column, top-left corner at (x,y)
\newcommand{\ffcol}[5]{\foreach \q in {1,...,#3}{\draw[#5] ({#1},{#2-\q*#4}) rectangle ++(#4,#4);}}
% \ffgrid{x}{ytop}{n}{per row}{cell}{style}: n flip-flops, filled row by row from the top-left
\newcommand{\ffgrid}[6]{\pgfmathtruncatemacro{\gl}{#3-1}%
  \foreach \b in {0,...,\gl}{\pgfmathtruncatemacro{\gc}{mod(\b,#4)}\pgfmathtruncatemacro{\gr}{div(\b,#4)}%
    \draw[#6] ({#1+\gc*#5},{#2-(\gr+1)*#5}) rectangle ++(#5,#5);}}

% every figure: same page size, title, footer
\newcommand{\canvas}{\useasboundingbox (-0.3,-0.2) rectangle (38.6,30.4);}
\newcommand{\figtitle}[3]{%
  \node[num, minimum size=0.9cm, font=\Large\bfseries] at (0.45,29.55) {#1};
  \node[anchor=west, font=\LARGE\bfseries, text=pcb] at (1.15,29.55) {#2};
  \node[anchor=west, font=\small, text=ink] at (0,28.7) {#3};}
\newcommand{\figfoot}[1]{\node[anchor=west, font=\scriptsize, text=muted] at (0,0.25) {#1};}
% \panel{x0}{y0}{x1}{y1}{title}
\newcommand{\panel}[5]{\draw[panel] (#1,#2) rectangle (#3,#4);
  \node[ptitle] at ({#1+0.3},{#4-0.25}) {#5};}

% ---------------- timing diagrams: 1 cycle = \cw cm, low at y, high at y+0.45 ----------------
\newcommand{\cw}{1.0}
\newcommand{\wclk}[3]{% x0 y ncycles
  \foreach \k in {1,...,#3}{\draw[pcb, thick] ({#1+(\k-1)*\cw},#2) -- ++(0,0.45) -- ++(0.5*\cw,0)
     -- ++(0,-0.45) -- ++(0.5*\cw,0);}}
\newcommand{\wsig}[3]{% x0 y {level per cycle, e.g. 0,1,1,0}
  \gdef\wprev{x}%
  \foreach \v [count=\k from 0] in {#3}{%
    \draw[pcb, thick] ({#1+\k*\cw},{#2+0.45*\v}) -- ++(\cw,0);
    \if x\wprev\else\ifnum\wprev=\v\else\draw[pcb, thick] ({#1+\k*\cw},#2) -- ++(0,0.45);\fi\fi
    \global\let\wprev\v}}
\newcommand{\wbus}[5]{% x0 y c0 c1 text: one value held from cycle c0 to cycle c1
  \draw[pcbmid, thick, fill=white] ({#1+(#3)*\cw},{#2+0.225}) -- ({#1+(#3)*\cw+0.08},{#2+0.45})
    -- ({#1+(#4)*\cw-0.08},{#2+0.45}) -- ({#1+(#4)*\cw},{#2+0.225}) -- ({#1+(#4)*\cw-0.08},#2)
    -- ({#1+(#3)*\cw+0.08},#2) -- cycle;
  \node[font=\tiny\ttfamily] at ({#1+((#3)+(#4))*\cw/2},{#2+0.225}) {#5};}
\newcommand{\widle}[4]{% x0 y c0 c1: bus idle (no value)
  \draw[muted, thick] ({#1+(#3)*\cw},{#2+0.225}) -- ({#1+(#4)*\cw},{#2+0.225});}
\newcommand{\wname}[3]{\node[anchor=east, font=\scriptsize\ttfamily] at ({#1-0.15},{#2+0.225}) {#3};}

\newcommand{\dff}[3]{% name x y : a D flip-flop box with a clock triangle
  \node[draw=pcb, fill=white, thick, minimum width=1.1cm, minimum height=1.2cm, font=\tiny, align=center] (#1) at (#2,#3) {D\hspace{0.35cm}Q\\[6pt]\ };
  \draw[pcb, thick] ($(#1.south)+(-0.15,0)$) -- ($(#1.south)+(0,0.2)$) -- ($(#1.south)+(0.15,0)$);}
\begin{document}
\begin{tikzpicture}[line join=round]
\canvas
\figtitle{9}{Clocks, resets and power: keeping every flip-flop ticking safely}
  {One crystal feeds PLLs that make every clock. Idle blocks get their clock stopped or their power cut.
   Signals that cross from one clock to another need synchronizers.}

% =====================================================================
% (A) PLL
% =====================================================================
\panel{0}{19.8}{19.0}{28.1}{PLL: from a 24 MHz crystal to a 3 GHz clock}
\node[hard, minimum width=2.0cm, minimum height=1.3cm, font=\tiny, text width=1.8cm] (xt) at (1.5,25.3) {\textbf{crystal}\\24 MHz};
\node[comb, minimum width=2.0cm, minimum height=1.3cm, font=\tiny, text width=1.8cm] (pfd) at (4.3,25.3) {\textbf{phase detector}\\who is early?};
\node[hard, minimum width=2.2cm, minimum height=1.3cm, font=\tiny, text width=2.0cm] (lf) at (7.2,25.3) {\textbf{charge pump + loop filter}\\$\rightarrow$ voltage};
\node[hard, minimum width=2.0cm, minimum height=1.3cm, font=\tiny, text width=1.8cm] (vco) at (10.1,25.3) {\textbf{VCO}\\voltage $\rightarrow$ frequency};
\node[comb, minimum width=2.0cm, minimum height=1.3cm, font=\tiny, text width=1.8cm] (od) at (12.9,25.3) {\textbf{output divider}\\$\div$P};
\draw[dat] (xt) -- (pfd); \draw[dat] (pfd) -- (lf); \draw[dat] (lf) -- (vco); \draw[dat] (vco) -- (od);
\node[comb, minimum width=1.6cm, minimum height=0.7cm, font=\tiny] (fb) at (7.2,23.4) {$\div$N = 125};
\draw[dat] (11.35,25.3) |- (fb.east); \draw[dat] (fb.west) -| (pfd.south);
\node[lab, text=ink, anchor=north] at (7.2,22.95) {feedback: the loop adjusts until 24 MHz $\times$ 125 = 3.0 GHz};
\node[anchor=north west, font=\scriptsize\bfseries] at (14.3,27.3) {example clocks (several PLLs):};
\node[note, font=\tiny, anchor=north west, text width=4.5cm] at (14.3,26.8) {%
  big CPU cores\quad 0.6--3.0 GHz (DVFS)\\little CPU cores\quad 0.3--2.0 GHz\\GPU\quad 0.3--1.0 GHz\\NPU\quad $\approx$1 GHz\\
  NoC\quad $\approx$1 GHz\\DRAM controller\quad 800 MHz\\peripherals\quad 24--100 MHz};
\draw[dat] (od.east) -- (14.2,25.3);
\node[note, anchor=north west, text width=18.2cm] at (0.3,22.4) {%
  $f_{out} = f_{ref} \times N / P$. The PLL also outputs \textbf{locked}: resets are released only after it is 1.
  The MAX 10 has PLLs too; your design uses the board's 50 MHz oscillator directly, with no PLL.};

% =====================================================================
% (B) clock tree + clock gating
% =====================================================================
\panel{19.3}{19.8}{38.3}{28.1}{Clock tree and clock gating}
% H-tree
\coordinate (rt) at (22.6,23.6);
\draw[pcb, very thick] (21.1,23.6) -- (24.1,23.6);
\foreach \x in {21.1,24.1}{
  \draw[pcb, thick] (\x,22.1) -- (\x,25.1);
  \foreach \y in {22.1,25.1}{
    \draw[pcb] ({\x-0.75},\y) -- ({\x+0.75},\y);
    \foreach \xx in {-0.75,0.75}{
      \draw[pcb] ({\x+\xx},{\y-0.6}) -- ({\x+\xx},{\y+0.6});
      \foreach \yy in {-0.6,0.6}{\draw[ff=cpud] ({\x+\xx-0.12},{\y+\yy-0.12}) rectangle ++(0.24,0.24);}}}}
\node[circle, fill=pcb, inner sep=1.5pt] at (rt) {};
\draw[dat] (22.6,27.0) -- (rt);
\node[lab, text=ink, anchor=south] at (22.6,27.0) {clock from PLL};
\node[lab, text=ink, anchor=north, text width=5.5cm, align=center] at (22.6,21.2) {H-tree: the same wire length to every flip-flop group $\Rightarrow$ small skew};
% ICG cell
\node[draw=pcb, fill=white, thick, minimum width=1.4cm, minimum height=1.0cm, font=\tiny, align=center] (lat) at (26.9,21.5) {latch\\open when\\CLK = 0};
\node[comb, minimum width=1.0cm, minimum height=1.0cm, font=\tiny\bfseries] (gand) at (28.9,21.5) {AND};
\draw[dat] (25.9,21.7) -- ($(lat.west)+(0,0.2)$); \node[lab, anchor=east] at (25.9,21.7) {EN};
\draw[dat] (lat) -- (gand);
\draw[dat] (25.9,20.5) -- (28.6,20.5) -- (28.6,21.0); \node[lab, anchor=east] at (25.9,20.5) {CLK};
\draw[dat] (gand.east) -- (30.0,21.5); \node[lab, anchor=west] at (30.0,21.5) {GCLK};
\node[lab, text=ink, anchor=south] at (27.9,22.05) {\textbf{ICG}: integrated clock-gating cell};
% waveform in half-cycles
\renewcommand{\cw}{0.45}
\wname{30.8}{26.4}{CLK}     \wsig{30.8}{26.4}{1,0,1,0,1,0,1,0,1,0,1,0,1,0,1,0}
\wname{30.8}{25.55}{EN}     \wsig{30.8}{25.55}{0,0,1,1,1,1,1,1,1,0,0,0,0,0,0,0}
\wname{30.8}{24.7}{latched} \wsig{30.8}{24.7}{0,0,0,1,1,1,1,1,1,0,0,0,0,0,0,0}
\wname{30.8}{23.85}{GCLK}   \wsig{30.8}{23.85}{0,0,0,0,1,0,1,0,1,0,0,0,0,0,0,0}
\node[note, font=\tiny, anchor=north west, text width=7.5cm] at (30.4,23.4) {%
  EN changes while CLK is high, but the latch passes it only while CLK is low, so GCLK gets 3 clean pulses and no glitch.
  A plain AND gate would cut a pulse short.\\[2pt]
  Stopped clock = no toggling = no dynamic power. FPGAs use a \textbf{clock enable} instead (your \rn{en} $\rightarrow$ LAB enable).};

% =====================================================================
% (C) DVFS
% =====================================================================
\panel{0}{10.4}{12.6}{19.5}{DVFS: lower voltage, lower frequency}
\node[note, anchor=north west, text width=12.0cm] at (0.3,18.75) {%
  Dynamic power $P = \alpha \cdot C \cdot V^2 \cdot f$ \quad ($\alpha$ = how often bits toggle, $C$ = wire + gate capacitance).
  Lower $f$ allows lower $V$, and $V$ counts twice.};
\node[anchor=north west, font=\scriptsize, inner sep=0pt] at (0.3,17.55) {%
  \renewcommand{\arraystretch}{1.1}%
  \begin{tabular}{@{}r r r r r@{}}
    \toprule
    \textbf{V} & \textbf{f (GHz)} & $\mathbf{V^2 f}$ & \textbf{power} & \textbf{speed}\\
    \midrule
    0.60 & 0.8 & 0.29 & 10\,\% & 27\,\%\\
    0.75 & 1.8 & 1.01 & 34\,\% & 60\,\%\\
    0.90 & 2.6 & 2.11 & 70\,\% & 87\,\%\\
    1.00 & 3.0 & 3.00 & 100\,\% & 100\,\%\\
    \bottomrule
  \end{tabular}};
\foreach \p/\s [count=\k from 0] in {10/27, 34/60, 70/87, 100/100}{
  \pgfmathsetmacro{\yy}{17.0-\k*0.75}
  \fill[pwrd!70] (7.6,\yy) rectangle ++({\p*0.045},0.28);
  \fill[cpud!60] (7.6,{\yy-0.3}) rectangle ++({\s*0.045},0.28);}
\node[lab, text=pwrd, font=\tiny\bfseries, anchor=west] at (7.6,17.55) {power};
\node[lab, text=cpud, font=\tiny\bfseries, anchor=west] at (9.0,17.55) {speed};
\node[note, anchor=north west, text width=12.0cm] at (0.3,13.85) {%
  From 3.0 to 1.8 GHz: 40\,\% slower, but 66\,\% less power, and energy per operation ($\propto V^2$) drops by 44\,\%.
  The PMU moves between these \textbf{operating points} many times a second. Example values; leakage power is ignored here.};

% =====================================================================
% (D) power domains
% =====================================================================
\panel{12.9}{10.4}{25.4}{19.5}{Power domains: switching a block off}
\draw[pwr] (13.2,15.2) rectangle (17.4,18.75);
\node[anchor=north west, font=\tiny\bfseries, text=pwrd] at (13.3,18.65) {always-on domain};
\node[note, font=\tiny, anchor=north west, text width=3.9cm] at (13.3,18.25) {PMU, RTC, wake-up logic. Never powered off.};
\node[pwr, fill=white, minimum width=3.8cm, minimum height=0.7cm, font=\tiny, anchor=north west] at (13.4,17.2) {PMU sequencer FSM};
\draw[npu] (19.6,15.2) rectangle (25.0,18.75);
\node[anchor=north west, font=\tiny\bfseries, text=npud] at (19.7,18.65) {NPU domain (switchable)};
\node[draw=pcb, fill=white, minimum width=1.6cm, minimum height=0.55cm, font=\tiny, anchor=north west] at (19.8,18.15) {power switch};
\node[lab, text=ink, anchor=west] at (21.5,17.88) {PMOS: VDD $\rightarrow$ VDD\_NPU};
\foreach \k in {0,1,2,3}{\draw[ff=npud] ({19.85+\k*0.4},16.75) rectangle ++(0.3,0.3);}
\node[lab, text=ink, anchor=west] at (21.5,16.9) {retention FFs keep state};
\node[comb, minimum width=1.5cm, minimum height=0.6cm, font=\tiny, anchor=north west] at (17.8,16.6) {isolation};
\draw[dat] (19.6,16.3) -- (19.3,16.3);
\draw[dat] (17.8,16.3) -- (17.4,16.3);
\node[lab, text=ink, anchor=north west, text width=2.0cm] at (17.75,15.9) {clamps outputs to 0 while off};
\node[comb, minimum width=1.5cm, minimum height=0.6cm, font=\tiny, anchor=north west] at (17.8,18.4) {level shift};
\node[note, font=\scriptsize, anchor=north west, text width=11.9cm] at (13.2,14.85) {%
  \textbf{Off}: 1 stop the clock, 2 save state into retention FFs, 3 isolation on, 4 power switch off.
  \textbf{On}: power on, wait until stable, restore, isolation off, release reset, start the clock.\\[3pt]
  Clock gating saves dynamic power; power gating also removes \textbf{leakage}, which grows as transistors shrink.
  Level shifters connect domains that run at different voltages.};

% =====================================================================
% (E) reset
% =====================================================================
\panel{25.7}{10.4}{38.3}{19.5}{Reset: assert at once, release on a clock edge}
\dff{r1}{28.2}{17.0} \dff{r2}{30.6}{17.0}
\draw[dat] (26.6,17.2) -- ($(r1.west)+(0,0.2)$); \node[lab, anchor=east] at (26.6,17.2) {1};
\draw[dat] ($(r1.east)+(0,0.2)$) -- ($(r2.west)+(0,0.2)$);
\draw[dat] ($(r2.east)+(0,0.2)$) -- (32.3,17.2); \node[lab, text=ink, anchor=west] at (32.3,17.2) {rst\_sync\_n $\rightarrow$ block};
\draw[ctl] (27.0,18.05) -- (30.6,18.05); \draw[ctl] (28.2,18.05) -- (r1.north); \draw[ctl] (30.6,18.05) -- (r2.north);
\node[lab, anchor=south west] at (26.0,18.1) {rst\_n (async clear)};
\node[note, font=\scriptsize, anchor=north west, text width=12.0cm] at (26.0,15.9) {%
  rst\_n = 0 clears both flip-flops at once (asynchronous). When rst\_n returns to 1, the 1 walks through the two
  flip-flops, so the block leaves reset exactly on a clock edge (synchronous release), even if rst\_n changed at a bad moment.\\[3pt]
  Your \rn{pe.sv} uses a fully synchronous reset (checked at the clock edge). Both work on FPGAs; ASIC teams usually
  pick asynchronous assert + synchronous release.\\[3pt]
  Order at power-on: supplies stable $\rightarrow$ PLL locked $\rightarrow$ release resets, block by block.};

% =====================================================================
% (F) CDC
% =====================================================================
\panel{0}{0.7}{38.3}{10.1}{Clock-domain crossing: 2-flip-flop synchronizer and async FIFO}
\node[anchor=north west, font=\scriptsize\bfseries] at (0.3,9.35) {1 bit: two flip-flops in a row (your \rn{sync\_2ff.sv})};
\dff{s1}{3.6}{7.6} \dff{s2}{6.0}{7.6}
\draw[dat] (1.2,7.8) -- ($(s1.west)+(0,0.2)$); \node[lab, anchor=south] at (1.6,7.8) {async in};
\draw[dat] ($(s1.east)+(0,0.2)$) -- ($(s2.west)+(0,0.2)$); \node[lab, anchor=south, text width=1.2cm, align=center] at (4.8,7.85) {may be metastable};
\draw[dat] ($(s2.east)+(0,0.2)$) -- (7.6,7.8); \node[lab, text=ink, anchor=west] at (7.6,7.8) {safe};
\draw[dat] (1.2,6.6) -- (6.0,6.6); \draw[dat] (3.6,6.6) -- (s1.south); \draw[dat] (6.0,6.6) -- (s2.south);
\node[lab, anchor=east] at (1.2,6.6) {clk};
\node[note, font=\scriptsize, anchor=north west, text width=9.4cm] at (0.3,6.2) {%
  If the input changes just at the clock edge, FF1 can hang between 0 and 1 for a short time (\textbf{metastability}).
  FF2 gives it a whole clock period to settle. The chance of a failure becomes tiny (MTBF of years).
  In your design the UART RX pin and the KEY buttons go through \rn{sync\_2ff}.};
\node[note, font=\scriptsize, anchor=north west, text width=9.4cm] at (0.3,3.3) {%
  \textbf{Never} send a multi-bit binary count through 2-FF synchronizers: the bits can arrive in different cycles
  (3 $\rightarrow$ 4 is 011 $\rightarrow$ 100: three bits change at once). Use Gray code, a handshake or an async FIFO.};
% async FIFO
\node[anchor=north west, font=\scriptsize\bfseries] at (10.2,9.35) {many bits: async FIFO};
\draw[cpu] (10.3,2.2) rectangle (15.6,8.8);
\node[anchor=north west, font=\tiny\bfseries, text=cpud] at (10.4,8.7) {write clock domain};
\node[cpu, fill=white, minimum width=4.9cm, minimum height=0.7cm, font=\tiny, anchor=north west] at (10.5,8.2) {write pointer (binary) $\rightarrow$ Gray};
\node[comb, minimum width=4.9cm, minimum height=0.7cm, font=\tiny, anchor=north west] at (10.5,7.2) {full? compare with synced read ptr};
\node[sram, minimum width=4.1cm, minimum height=2.6cm, font=\tiny, anchor=north west, text width=3.9cm] at (16.0,8.2)
  {\textbf{dual-port RAM}\\write port on wclk\\read port on rclk\\(an M9K can do this)};
\draw[gpu] (20.5,2.2) rectangle (25.8,8.8);
\node[anchor=north west, font=\tiny\bfseries, text=gpud] at (20.6,8.7) {read clock domain};
\node[gpu, fill=white, minimum width=4.9cm, minimum height=0.7cm, font=\tiny, anchor=north west] at (20.7,8.2) {read pointer (binary) $\rightarrow$ Gray};
\node[comb, minimum width=4.9cm, minimum height=0.7cm, font=\tiny, anchor=north west] at (20.7,7.2) {empty? compare with synced write ptr};
\foreach \k in {0,1}{\draw[ff=cpud] ({17.2+\k*0.45},4.6) rectangle ++(0.35,0.35);}
\foreach \k in {0,1}{\draw[ff=gpud] ({18.6+\k*0.45},4.0) rectangle ++(0.35,0.35);}
\draw[dat] (15.6,4.25) -- (18.6,4.25); \draw[dat] (19.45,4.17) -- (20.5,4.17);
\draw[dat] (20.5,4.77) -- (18.1,4.77); \draw[dat] (17.2,4.77) -- (15.6,4.77);
\node[lab, text=ink, anchor=north, text width=4.0cm, align=center] at (18.05,3.9) {Gray pointers cross through 2-FF synchronizers};
\node[note, font=\scriptsize, anchor=north west, text width=5.0cm] at (10.4,6.2) {Only 1 bit of a Gray pointer changes per step, so a half-updated value is still either the old or the new pointer.};
% gray table
\node[anchor=north west, font=\scriptsize, inner sep=0pt] at (26.4,9.3) {%
  \renewcommand{\arraystretch}{1.0}%
  \begin{tabular}{@{}c c c@{}}
    \toprule
    \textbf{n} & \textbf{binary} & \textbf{Gray}\\
    \midrule
    0 & 000 & 000\\ 1 & 001 & 001\\ 2 & 010 & 011\\ 3 & 011 & 010\\
    4 & 100 & 110\\ 5 & 101 & 111\\ 6 & 110 & 101\\ 7 & 111 & 100\\
    \bottomrule
  \end{tabular}};
\node[note, font=\scriptsize, anchor=north west, text width=7.6cm] at (30.4,9.25) {%
  \textbf{Gray = binary XOR (binary $\gg$ 1)}: one XOR gate per bit.\\[3pt]
  Each step changes exactly one bit (look down the right column), also from 7 back to 0.\\[4pt]
  \textbf{Your design} has one clock (50 MHz) for everything, so the only crossings are the asynchronous input pins.
  A real SoC has dozens of clock domains and hundreds of crossings; CDC checking tools verify every one.};

\figfoot{PLL settings, operating points and clock rates are example values. Power-domain cell names follow common low-power design practice.}
\end{tikzpicture}
\end{document}
